% ============================================================================
% RELATED: HYPERDIMENSIONAL COMPUTING AND VECTOR SYMBOLIC ARCHITECTURES
% ============================================================================
\subsection{Hyperdimensional Computing and Vector Symbolic Architectures}
\label{sec:related_hdc_vsa}

The Value representation is situated within the framework of
\emph{Hyperdimensional Computing} (HDC) and \emph{Vector Symbolic
Architectures} (VSA).  In the canonical VSA formulation
\citep{kanerva2009hyperdimensional}, data is represented as
high-dimensional random vectors (typically $D \ge 10{,}000$) and
manipulated through three algebraic operations: \emph{binding}
(element-wise multiplication or circular convolution), \emph{bundling}
(element-wise addition or disjunction), and \emph{permutation}
(coordinate rotation).  These operations form an approximate algebra
over a vector space, enabling compositional representations of
structured data without learned parameters.

The Value maps onto this framework with specific design choices.  At
8192 bits, each frame operates in a dimensionality approaching the
classical HD regime ($D \ge 4{,}000$), gaining the
concentration-of-measure benefits that make high-dimensional binary
vectors approximately orthogonal under random sampling.  The boolean
instruction set provides all sixteen two-input operations, subsuming the
standard VSA binding and bundling operators as special cases: bitwise AND
serves as a similarity probe, bitwise OR implements bundling through
disjunction, and circular shift (achievable through register-offset
addressing) provides positional binding.

Gayler's analysis \citep{gayler2003vector} demonstrated that VSAs
support systematic compositional generalization because binding and
bundling preserve algebraic structure.  The Value architecture inherits
this property: the encoding of input bytes into the data region is
deterministic, and the sixteen boolean operations are closed over the
binary frame, meaning that any sequence of operations on valid frames
produces a valid frame.

A notable difference from classical HDC is the absence of a fixed
bundling threshold.  Standard VSAs operate in a flat vector space where
repeated bundling via disjunction monotonically increases the population
count (the number of active bits), eventually degrading all similarity
queries when the count exceeds approximately 50\% of the dimensionality.
The Value architecture addresses this through the multi-tier Region
structure: rather than accumulating all information into a single vector,
the system distributes content across a population of frames that
interact combinatorially.  Saturation is a property of individual frames,
not of the system as a whole.
